\begin{remarks}[7.3.5]
\label{IV.2.s7.3.5-fr}
\begin{enumerate}
\item[(i)] Las condiciones $(\mathrm P_{\mathrm I})$ y
$(\mathrm P_{\mathrm{III}})$ implican que \emph{todo morfismo regular
es un $\boldsymbol P$-morfismo}.
\item[(ii)] Observemos que las hipótesis de
$(\mathrm P_{\mathrm I})$ (resp. de $(\mathrm P_{\mathrm{II}})$)
implican que $h=g\circ f$ (resp. $g$) es plano (2.2.13); por otra
parte, las hipótesis de $(\mathrm P_{\mathrm I})$ o de
$(\mathrm P_{\mathrm{II}})$ implican que, para todo $z\in Z$,
$f_z:h^{-1}(z)\to g^{-1}(z)$ es plano (2.1.4). Como, para todo
$y\in g^{-1}(z)$, $f_z^{-1}(y)$ es isomorfo a $f^{-1}(y)$
(\textbf I, 3.6.4), esto muestra que, para verificar las condiciones
$(\mathrm P_{\mathrm I})$ y $(\mathrm P_{\mathrm{II}})$, basta
hacerlo cuando $Z$ es el \emph{espectro de un cuerpo}.
\item[(iii)] En ciertos casos, la propiedad $\boldsymbol Q$
será tal que se cumpla la condición siguiente:
\begin{enumerate}
\item[$(\mathrm P'_{\mathrm I})$.] --- Si $f:X\to Y$ y $g:Y\to Z$
son dos $\boldsymbol P$-morfismos, entonces $g\circ f$ es un
$\boldsymbol P$-morfismo.
\end{enumerate}
\end{enumerate}
\end{remarks}
